\documentclass[aspectratio=169,xcolor=table]{beamer}

\usepackage[utf8]{inputenc}
\usepackage[T1]{fontenc}
\usepackage[portuguese]{babel}
\usepackage{lmodern}
\usepackage{booktabs}
\usepackage{array}
\usepackage{tikz}
\usepackage{pgfplots}
\usepackage{amsmath}
\usepackage{csquotes}
\usepackage{xcolor}

\pgfplotsset{compat=1.18}
\usetheme{DCC}

\setbeamertemplate{itemize items}[circle]
\setbeamertemplate{itemize subitem}[circle]
\setlength{\itemsep}{0.45em}
\setlength{\parskip}{0.35em}
\graphicspath{{imgs/}{./imgs/}}

\newcommand{\metricup}{\ensuremath{\uparrow}}
\newcommand{\metricdown}{\ensuremath{\downarrow}}
\newcommand{\uf}[1]{\textsc{#1}}
\newcommand{\smallnote}[1]{\vspace{0.2em}{\scriptsize\color{black!70}#1}}

\author[Antoniel Magalh\~aes]{\textbf{Antoniel Magalh\~aes}}
\title{Agrupamento n\~ao supervisionado de dengue no Nordeste}
\subtitle{K-means aplicado a notifica\c{c}\~oes SINAN por munic\'ipio e semana epidemiol\'ogica}
\institute{Intelig\^encia Artificial IV}
\date{\today}

\begin{document}

\begin{frame}[plain,noframenumbering]
    \titlepage
\end{frame}

\begin{frame}{Agenda}
    \tableofcontents
\end{frame}

%=================================================
\section{Problema e dados}
%=================================================
\sectionframe

\begin{frame}{Motiva\c{c}\~ao}
    \begin{block}{Pergunta do trabalho}
        Como identificar padr\~oes de intensidade de dengue entre munic\'ipios sem r\'otulos pr\'evios de risco?
    \end{block}

    \begin{itemize}
        \item Dengue tem forte varia\c{c}\~ao \textbf{espacial} e \textbf{temporal}.
        \item Totais anuais escondem picos semanais e deslocamentos regionais.
        \item O objetivo \textbf{n\~ao} \`e prever casos: \`e descobrir grupos de munic\'ipios com intensidade semelhante.
    \end{itemize}
\end{frame}

\begin{frame}{Base de dados}
    \begin{columns}[T]
        \begin{column}{0.50\textwidth}
            \textbf{Fonte principal}
            \begin{itemize}
                \item SINAN / OpenDataSUS
                \item Notifica\c{c}\~oes de arboviroses -- dengue
                \item CSVs nacionais: \texttt{DENGBR16.csv} a \texttt{DENGBR26.csv}
            \end{itemize}
        \end{column}
        \begin{column}{0.45\textwidth}
            \textbf{Recorte usado}
            \begin{itemize}
                \item Nordeste: 9 UFs
                \item Ano principal da an\'alise: 2024
                \item Unidade anal\'itica: munic\'ipio $\times$ semana epidemiol\'ogica
            \end{itemize}
        \end{column}
    \end{columns}

    \vspace{0.5em}
    \begin{block}{Campos centrais}
        \texttt{SG\_UF\_NOT} para UF, \texttt{ID\_MUNICIP} para munic\'ipio e \texttt{SEM\_NOT} para semana epidemiol\'ogica.
    \end{block}
\end{frame}

\begin{frame}{Pipeline do projeto}
    \begin{center}
    \begin{tikzpicture}[node distance=1.55cm, every node/.style={font=\small}]
        \tikzset{flowstep/.style={rectangle, rounded corners, draw=black!60, fill=MedianUltraLightOrange, minimum width=2.2cm, minimum height=0.8cm, align=center}}
        \node[flowstep] (raw) {SINAN\\CSV nacional};
        \node[flowstep, right of=raw, xshift=1.8cm] (filter) {Filtro\\por UF};
        \node[flowstep, right of=filter, xshift=1.8cm] (agg) {Agrega\c{c}\~ao\\munic\'ipio/semana};
        \node[flowstep, right of=agg, xshift=1.8cm] (km) {K-means\\semanal};
        \node[flowstep, right of=km, xshift=1.8cm] (map) {Mapa\\interativo};
        \draw[->, thick] (raw) -- (filter);
        \draw[->, thick] (filter) -- (agg);
        \draw[->, thick] (agg) -- (km);
        \draw[->, thick] (km) -- (map);
    \end{tikzpicture}
    \end{center}

    \begin{itemize}
        \item Malhas municipais do IBGE s\~ao usadas para visualiza\c{c}\~ao espacial.
        \item Processamento e clusters ficam em cache para evitar recomputa\c{c}\~ao.
    \end{itemize}
\end{frame}

%=================================================
\section{Modelo n\~ao supervisionado}
%=================================================
\sectionframe

\begin{frame}{Por que aprendizado n\~ao supervisionado?}
    \begin{itemize}
        \item N\~ao existe no dado um r\'otulo pronto como ``baixo'', ``alto'' ou ``cr\'itico''.
        \item Queremos encontrar grupos a partir da pr\'opria distribui\c{c}\~ao dos casos.
        \item O K-means transforma contagens brutas em faixas interpret\'aveis no mapa.
    \end{itemize}

    \begin{block}{Interpreta\c{c}\~ao desejada}
        Grupos ordinais de intensidade: menor intensidade $\rightarrow$ maior intensidade.
    \end{block}
\end{frame}

\begin{frame}{Como o K-means foi aplicado}
    Para cada UF, ano e semana epidemiol\'ogica:

    \begin{enumerate}
        \item Montar vetor de casos por munic\'ipio.
        \item Filtrar munic\'ipios eleg\'iveis: total anual $\geq 30$ casos.
        \item Transformar a contagem: $X = \log(1 + casos)$.
        \item Aplicar K-means em uma dimens\~ao.
        \item Reordenar clusters pela m\'edia de casos.
    \end{enumerate}

    \begin{block}{Por que usar \(\log(1 + casos)\)?}
        A distribui\c{c}\~ao de dengue \`e muito assim\'etrica: poucos munic\'ipios concentram muitos casos. O log reduz o peso dos outliers.
    \end{block}
\end{frame}

\begin{frame}{O problema de escolher K}
    \begin{columns}[T]
        \begin{column}{0.48\textwidth}
            \textbf{K baixo}
            \begin{itemize}
                \item Mais est\'avel
                \item Mais simples
                \item Pode esconder padr\~oes relevantes
            \end{itemize}
        \end{column}
        \begin{column}{0.48\textwidth}
            \textbf{K alto}
            \begin{itemize}
                \item Mais granular
                \item Melhora algumas m\'etricas internas
                \item Pode gerar clusters pequenos e inst\'aveis
            \end{itemize}
        \end{column}
    \end{columns}

    \vspace{0.5em}
    \begin{alertblock}{Crit\'erio pr\'atico}
        O melhor K inicial precisa equilibrar qualidade estat\'istica, estabilidade temporal e leitura visual no mapa.
    \end{alertblock}
\end{frame}

%=================================================
\section{Valida\c{c}\~ao de K}
%=================================================
\sectionframe

\begin{frame}{M\'etricas usadas}
    \small
    \begin{tabular}{p{0.25\textwidth} p{0.17\textwidth} p{0.48\textwidth}}
        \toprule
        \textbf{M\'etrica} & \textbf{Dire\c{c}\~ao} & \textbf{O que avalia} \\
        \midrule
        Silhouette & maior melhor & Coes\~ao interna e separa\c{c}\~ao entre clusters \\
        Davies--Bouldin & menor melhor & Dispers\~ao intra-cluster vs dist\^ancia entre clusters \\
        Calinski--Harabasz & maior melhor & Vari\^ancia entre clusters vs dentro dos clusters \\
        Inertia / cotovelo & joelho da curva & Ponto de ganho marginal decrescente \\
        ARI consecutivo & maior melhor & Estabilidade entre semanas consecutivas \\
        \bottomrule
    \end{tabular}

    \smallnote{Como o dado \`e 1D, silhouette e Davies--Bouldin tendem a favorecer K alto. Por isso o cotovelo e a interpretabilidade s\~ao decisivos.}
\end{frame}

\begin{frame}{Cobertura analisada em 2024}
    \small
    \begin{columns}[T]
        \begin{column}{0.42\textwidth}
            \begin{tabular}{lrr}
                \toprule
                \textbf{UF} & \textbf{Eleg\'iveis} & \textbf{Total} \\
                \midrule
                BA & 352 & 415 \\
                PE & 88 & 180 \\
                MA & 71 & 198 \\
                RN & 54 & 157 \\
                PB & 53 & 204 \\
                PI & 52 & 170 \\
                CE & 47 & 178 \\
                AL & 38 & 93 \\
                SE & 9 & 56 \\
                \bottomrule
            \end{tabular}
        \end{column}
        \begin{column}{0.55\textwidth}
            \begin{tikzpicture}
            \begin{axis}[
                ybar,
                width=\textwidth,
                height=0.62\textheight,
                symbolic x coords={BA,PE,MA,RN,PB,PI,CE,AL,SE},
                xtick=data,
                ylabel={Munic\'ipios eleg\'iveis},
                ymin=0,
                bar width=9pt,
                nodes near coords,
                nodes near coords style={font=\tiny},
                tick label style={font=\scriptsize},
                ymajorgrids=true,
                grid style={black!10}
            ]
            \addplot[fill=MedianOrange] coordinates {(BA,352) (PE,88) (MA,71) (RN,54) (PB,53) (PI,52) (CE,47) (AL,38) (SE,9)};
            \end{axis}
            \end{tikzpicture}
        \end{column}
    \end{columns}
\end{frame}

\begin{frame}{Resultados agregados: Nordeste 2024}
    \small
    \begin{tabular}{rrrrrr}
        \toprule
        \textbf{K} & \textbf{Silhouette} & \textbf{DB} & \textbf{Inertia} & \textbf{ARI} & \textbf{menor cluster} \\
        \midrule
        2 & 0{,}744 & 0{,}431 & 25{,}9 & \textbf{0{,}451} & 19{,}7 \\
        3 & 0{,}774 & 0{,}366 & 11{,}3 & 0{,}301 & 11{,}4 \\
        \rowcolor{MedianUltraLightOrange}
        \textbf{4} & \textbf{0{,}801} & \textbf{0{,}298} & \textbf{5{,}6} & 0{,}285 & \textbf{7{,}6} \\
        5 & 0{,}822 & 0{,}247 & 3{,}2 & 0{,}280 & 5{,}4 \\
        6 & 0{,}837 & 0{,}206 & 1{,}9 & 0{,}279 & 4{,}1 \\
        7 & 0{,}846 & 0{,}173 & 1{,}3 & 0{,}277 & 3{,}2 \\
        8 & \textbf{0{,}854} & \textbf{0{,}149} & 0{,}9 & 0{,}276 & 2{,}5 \\
        \bottomrule
    \end{tabular}

    \vspace{0.4em}
    \smallnote{Silhouette e DB melhoram com K alto, mas o tamanho do menor cluster cai. ARI favorece K=2. O cotovelo da inertia indica o ponto de equil\'ibrio.}
\end{frame}

\begin{frame}{Cotovelo da inertia}
    \begin{columns}[T]
        \begin{column}{0.58\textwidth}
            \begin{tikzpicture}
            \begin{axis}[
                width=\textwidth,
                height=0.68\textheight,
                xlabel={K},
                ylabel={Inertia m\'edia},
                xtick={2,3,4,5,6,7,8},
                ymin=0,
                ymajorgrids=true,
                grid style={black!10},
                legend pos=north east
            ]
            \addplot[color=MedianOrange, mark=* , thick] coordinates {(2,25.9) (3,11.3) (4,5.6) (5,3.2) (6,1.9) (7,1.3) (8,0.9)};
            \addlegendentry{Inertia}
            \addplot[color=DCCRedHard, dashed, thick] coordinates {(4,0) (4,26)};
            \addlegendentry{cotovelo: K=4}
            \end{axis}
            \end{tikzpicture}
        \end{column}
        \begin{column}{0.38\textwidth}
            \begin{block}{Resultado}
                O cotovelo escolheu \textbf{K=4 em 9/9 UFs}.
            \end{block}
            \begin{itemize}
                \item K=2 ainda reduz muito a complexidade.
                \item Depois de K=4, o ganho marginal cai.
                \item K alto fragmenta os grupos.
            \end{itemize}
        \end{column}
    \end{columns}
\end{frame}

\begin{frame}{Votos por m\'etrica}
    \small
    \begin{tabular}{ll}
        \toprule
        \textbf{M\'etrica} & \textbf{Vencedor por UF} \\
        \midrule
        Cotovelo / inertia & \textbf{K=4 em 9/9 UFs} \\
        ARI / estabilidade & K=2 em 9/9 UFs \\
        Silhouette & K=8 em 8/9 UFs; K=4 em 1/9 \\
        Davies--Bouldin & K=8 em 9/9 UFs \\
        Calinski--Harabasz & K=8 em 5/9; K=7 em 2/9; disperso nos demais \\
        \bottomrule
    \end{tabular}

    \vspace{0.6em}
    \begin{alertblock}{Leitura}
        N\~ao existe uma \u201cm\'etrica m\'agica\u201d. Em clustering 1D, m\'etricas internas puras empurram K para cima. O default precisa ser uma decis\~ao de compromisso.
    \end{alertblock}
\end{frame}

\begin{frame}{Bahia como refer\^encia}
    \begin{columns}[T]
        \begin{column}{0.48\textwidth}
            \small
            \begin{tabular}{rrrrr}
                \toprule
                \textbf{K} & \textbf{Sil.} & \textbf{DB} & \textbf{ARI} & \textbf{menor cluster} \\
                \midrule
                2 & 0{,}734 & 0{,}501 & \textbf{0{,}481} & 87 \\
                3 & 0{,}765 & 0{,}480 & 0{,}342 & 42 \\
                \rowcolor{MedianUltraLightOrange}
                \textbf{4} & \textbf{0{,}789} & \textbf{0{,}446} & 0{,}331 & \textbf{20} \\
                5 & 0{,}810 & 0{,}409 & 0{,}322 & 13 \\
                8 & 0{,}846 & 0{,}341 & 0{,}314 & 4 \\
                \bottomrule
            \end{tabular}
        \end{column}
        \begin{column}{0.48\textwidth}
            \begin{itemize}
                \item BA tem 352 munic\'ipios eleg\'iveis: maior base da an\'alise.
                \item K=8 melhora silhouette, mas cria clusters m\'edios m\'inimos de apenas 4 munic\'ipios.
                \item K=4 preserva granularidade sem quebrar demais o mapa.
            \end{itemize}
        \end{column}
    \end{columns}
\end{frame}

%=================================================
\section{Decis\~ao e limita\c{c}\~oes}
%=================================================
\sectionframe

\begin{frame}{Racional final: por que K=4?}
    \begin{enumerate}
        \item \textbf{Cotovelo robusto:} K=4 foi escolhido em 9/9 UFs.
        \item \textbf{Interpretabilidade:} quatro grupos comunicam bem baixa, moderada, alta e cr\'itica intensidade.
        \item \textbf{Evita simplifica\c{c}\~ao excessiva:} K=2 \`e est\'avel, mas pouco informativo.
        \item \textbf{Evita fragmenta\c{c}\~ao:} K=8 melhora m\'etricas internas, mas cria grupos pequenos.
        \item \textbf{Ader\^encia ao dashboard:} K=4 funciona bem para mapa coropl\'etico e legenda ordinal.
    \end{enumerate}

    \begin{block}{Conclus\~ao}
        K=4 n\~ao \u00e9 o \u201c\'otimo absoluto\u201d de uma m\'etrica isolada; \u00e9 o melhor \textbf{default inicial} para an\'alise e apresenta\c{c}\~ao.
    \end{block}
\end{frame}

\begin{frame}{Limita\c{c}\~oes}
    \begin{itemize}
        \item Usa contagem absoluta de casos, ainda sem normaliza\c{c}\~ao por popula\c{c}\~ao.
        \item Clustering \`e semanal e independente; fronteiras mudam a cada semana.
        \item A vari\'avel de entrada \u00e9 1D: apenas casos semanais transformados por log.
        \item Munic\'ipios com menos de 30 casos/ano recebem tratamento especial.
        \item M\'etricas internas podem favorecer K alto em dados orden\'aveis.
        \item Resultado principal foi calibrado em 2024; outros anos podem ter din\^amicas diferentes.
    \end{itemize}
\end{frame}

\begin{frame}{Pr\'oximos passos}
    \begin{itemize}
        \item Repetir a an\'alise para 2022 e 2023.
        \item Comparar \texttt{min\_casos=10} vs \texttt{min\_casos=30}.
        \item Normalizar por popula\c{c}\~ao: casos por 100 mil habitantes.
        \item Medir estabilidade em janelas m\'oveis de 4 semanas.
        \item Avaliar K adaptativo para UFs pequenas, como Sergipe.
        \item Cruzar clusters com a\c{c}\~oes de controle vetorial e indicadores externos.
    \end{itemize}
\end{frame}

\begin{frame}{Conclus\~ao}
    \begin{block}{Mensagem principal}
        A IA n\~ao supervisionada transformou notifica\c{c}\~oes brutas em grupos interpret\'aveis de intensidade semanal de dengue.
    \end{block}

    \begin{itemize}
        \item K-means ajuda a revelar padr\~oes espaciais e temporais sem r\'otulos pr\'evios.
        \item K=4 \u00e9 uma escolha empiricamente sustentada pelo cotovelo e adequada para comunica\c{c}\~ao visual.
        \item O modelo ainda precisa evoluir para incid\^encia populacional e valida\c{c}\~ao externa.
    \end{itemize}
\end{frame}

\begin{frame}{Perguntas}
    \begin{center}
        \Huge Obrigado!

        \vspace{0.8em}
        \large Perguntas e coment\'arios
    \end{center}
\end{frame}

\end{document}
